% Revize 2026-11-01: vsechny vyrazy prevzaty z drivejsich hodin 73 a 74,
% presunuty dle obtiznosti. Inspirace Realisticky.cz 020203--020204,
% konkretni vyzvy prevzaty z ucitelem dodanych fotografii (c. 1a,1b).
\documentclass[11pt,a4paper]{article}
\input{../../sablony/pripravy/priprava.sty}
\nastavPripravu{1. ročník}{Rovnice a nerovnice}{74}{Lineární rovnice: složené závorky}
\begin{document}
\blok{Návaznost: jedna závorka nestačí}
V hodině 73 žáci řešili rovnice s jednoduchými závorkami a rozlišovali jeden kořen, žádné řešení a nekonečně mnoho řešení. Nyní přidáme \emph{více úrovní závorek}. Algebraická pravidla už znají z úprav výrazů: nejdříve upravíme nejvnitřnější závorku, teprve potom postupujeme ven. Ekvivalentní úpravy rovnic zůstávají stejné.
\blok{Společný postup}
\begin{enumerate}[label=\arabic*.]
\item Přesně vyznačit pořadí závorek $(\cdots)$, $[\cdots]$, $\{\cdots\}$.
\item Roznásobovat od nejvnitřnější závorky; kontrolovat záporné znaménko před závorkou.
\item Vzniklou lineární rovnici vyřešit a provést zkoušku v původním zadání.
\end{enumerate}
\begin{prikladReseniBox}[title={Příklad 1}]
Vyřešte $4\{2x-[3-(x-2)]\}=12$.
\tcblower
\[
\begin{aligned}
4\{2x-[3-x+2]\}&=12\\
4\{2x-(5-x)\}&=12\\
4(3x-5)&=12\\
12x-20&=12\\
12x&=32\\
x&=\frac83.
\end{aligned}
\]
Zkouška je nejpřehlednější v ekvivalentním tvaru: $12\cdot\frac83-20=12$.
\[
\begin{aligned}
K &= \{\frac83\}
\end{aligned}
\]
\end{prikladReseniBox}
\begin{prikladReseniBox}[title={Příklad 2}]
Vyřešte $2\{3[x-2(1-x)]+1\}=26$.
\tcblower
\[
\begin{aligned}
2\{3[x-2+2x]+1\}&=26\\
2\{3(3x-2)+1\}&=26\\
2(9x-5)&=26\\
18x-10&=26\\
18x&=36\\
x&=2.
\end{aligned}
\]
Zkouška v původním zadání: $2\{3[2-2(1-2)]+1\}=2(12+1)=26$.
\[
\begin{aligned}
K &= \{2\}
\end{aligned}
\]
\end{prikladReseniBox}
\poznamka{Nezvyšuj současně počet závorek a náročnost aritmetiky. V prvních dvou příkladech průběžně kontroluj distribuci znamének; v navazujících úlohách už nech prostor pro samostatný postup.}
\blok{Samostatné procvičování – výběr podle tempa}
\begin{zadaniReseniBox}[title={Samostatná práce 1}]
\begin{enumerate}[label=\alph*)]
\item $4[2x-(x-3)]=20$.
\item $4[3-(2x-5)]=20$.
\item $3\{2[3x-(x+1)]-4\}=12$.
\end{enumerate}
\tcblower
\begin{enumerate}[label=\alph*)]
\item
\[
\begin{aligned}4(x+3)&=20\\x+3&=5\\x&=2. \\[2mm]
K &= \{2\}
\end{aligned}
\]
\item
\[
\begin{aligned}4(8-2x)&=20\\32-8x&=20\\-8x&=-12\\x&=\tfrac32. \\[2mm]
K &= \{\tfrac32\}
\end{aligned}
\]
\item
\[
\begin{aligned}3\{2(2x-1)-4\}&=12\\3(4x-6)&=12\\12x-18&=12\\x&=\tfrac52. \\[2mm]
K &= \{\tfrac52\}
\end{aligned}
\]
\end{enumerate}
\end{zadaniReseniBox}
\blok{Volitelná rezerva – náročný příklad ze sbírky}
\begin{prikladReseniBox}[title={Rozšíření 1}]
Vyřešte $5\{5[5(5x-4)-4]-4\}=5$.
\tcblower
\[
\begin{aligned}
5\{5[25x-20-4]-4\}&=5\\
5\{5(25x-24)-4\}&=5\\
5(125x-124)&=5\\
625x-620&=5\\
625x&=625\\
x&=1.
\end{aligned}
\]
Přímá kontrola: $5\{5[5(5-4)-4]-4\}=5\{5(1)-4\}=5$.
\[
\begin{aligned}
K &= \{1\}
\end{aligned}
\]
\end{prikladReseniBox}
\blok{Domácí úkol}
\begin{zadaniReseniBox}[title={Domácí úkol}]
\begin{enumerate}[label=\alph*)]
\item
$5\{2[x-(3-x)]+1\}=35$.
\item
$3[2(3x-6)-2(4x-5)+1]-3=6[3-8(x-3)]$ (úloha č. 1b z fotografie).
\item
$7\{x-2[5+x-8(x+3)-1]+2\}=3$.
\item
$4[2x-3(x-1)]=20$.
\item
$2\{3[2x-(x-4)]-5\}=2x+8$.
\item
$5\{2[3-(x-2)]-3(x+1)\}=10$.
\end{enumerate}
\tcblower
\begin{enumerate}[label=\alph*)]
\item
\[
\begin{aligned}5(4x-5)&=35\\4x-5&=7\\x&=3. \\[2mm]
K &= \{3\}
\end{aligned}
\]
\item
\[
\begin{aligned}3[-2x-1]-3&=6[27-8x]\\-6x-6&=162-48x\\42x&=168\\x&=4. \\[2mm]
K &= \{4\}
\end{aligned}
\]
\item
\[
\begin{aligned}7\{x-2[-20-7x]+2\}&=3\\7(15x+42)&=3\\105x+294&=3\\105x&=-291\\x&=-\tfrac{97}{35}. \\[2mm]
K &= \{-\tfrac{97}{35}\}
\end{aligned}
\]
\item
\[
\begin{aligned}
4(3-x)&=20\\12-4x&=20\\x&=-2 \\[2mm]
K &= \{-2\}
\end{aligned}
\]
\item
\[
\begin{aligned}
2(3x+7)&=2x+8\\6x+14&=2x+8\\4x&=-6\\x&=-\frac32 \\[2mm]
K &= \{-\frac32\}
\end{aligned}
\]
\item
\[
\begin{aligned}
5(7-5x)&=10\\35-25x&=10\\-25x&=-25\\x&=1 \\[2mm]
K &= \{1\}
\end{aligned}
\]
\end{enumerate}
\end{zadaniReseniBox}
\poznamka{45 minut: návaznost a technika 5 min, příklad 1 9 min, příklad 2 10 min, samostatně 15 min, rozbor častých chyb 4 min, organizace 2 min. Náročný příklad ze sbírky je volitelný; hlavně nesmí vytlačit samostatnou práci. V 75. hodině se přidají zlomky a kombinované úpravy.}
\end{document}
